\chapter{Posterior Targets and HMC Requirements}
\label{ch:bf-hmc-target-requirements}

HMC samples an unconstrained or transformed parameter vector by integrating
Hamiltonian dynamics for a log target.  In BayesFilter, that log target usually
contains a filtering likelihood.  The target contract must therefore be stated
before discussing sampler tuning.

\begin{definition}[Filtering posterior target]
\label{def:bf-filtering-target}
Let $u \in \R^{d_u}$ denote the coordinates used by the sampler and let
$\theta = T(u)$ be the model parameter transform.  A filtering posterior target
has the form
\[
  \log \pi(u \mid y_{1:T})
  =
  \log p(T(u))
  + \ell(T(u); y_{1:T})
  + \log\left|\det DT(u)\right|
  + \text{constant},
\]
where $\ell$ is the log likelihood returned by the chosen filtering backend.
\end{definition}

For HMC, the value alone is not enough.  The backend and transform jointly
define the gradient seen by the sampler.  If a custom gradient is used, the exact HMC target remains the same only when the Metropolis correction evaluates that same scalar value. A deterministic nearby force field can still leave the invariant law exact under exact endpoint Metropolis correction, but then the approximation changes leapfrog efficiency rather than the target itself. If the acceptance step also uses a surrogate value, or if the dynamics are left unadjusted, the target changes unless the sampler is explicitly corrected for that surrogate.  In
particular, if the filtering backend exposes likelihood derivatives in model
parameter coordinates \(\theta\), the HMC wrapper must still construct the full
sampler-coordinate target in \(u\)-space by adding the prior and transform
Jacobian terms from Definition~\ref{def:bf-filtering-target}.

\section{Finite Failure Policy}
\label{sec:bf-finite-failure-policy}

Invalid parameters are ordinary events during HMC warmup and exploration.  A
BayesFilter target must decide how to handle them.  BayesFilter must state whether an invalid region is handled by the exact support-rejection contract or by a declared modified target. Under the exact contract, the value path returns $-\infty$ (equivalently $U_F=+\infty$) and the proposal is rejected. Under the modified-target contract, the value path returns a deterministic fallback scalar $c$ together with the exact derivative of that same returned scalar, so the sampled law is proportional to $\exp(\widetilde\ell)$ where $\widetilde\ell=\ell$ on the valid region and $\widetilde\ell=c$ on the declared invalid region. If $c$ is constant on that region, the corresponding gradient must be zero. This finite fallback can be useful for compiled finite-valued control flow, but it is not the same target as support rejection, and on an unbounded invalid region it can be improper. A linear-algebra failure that aborts the process is not an HMC-compatible failure mode.

The policy must be explicit for:
\begin{itemize}
  \item parameter transforms and prior support;
  \item stationarity or determinacy restrictions;
  \item covariance positive-definiteness checks;
  \item failed factorizations or solve residuals;
  \item non-finite likelihood or gradient values.
\end{itemize}

\section{Compilation and Differentiability}
\label{sec:bf-hmc-target-compilation}

When a backend is advertised as JIT or XLA compatible, the claim refers to the
complete value-and-gradient path needed by the sampler.  Partial tracing of an
outer function is not enough if the filter, derivative provider, or failure
branch executes outside the compiled region.

The derivative policy should be one of:
\begin{itemize}
  \item analytic score or Hessian recursion;
  \item custom gradient wrapping a certified analytic derivative;
  \item autodiff through smooth primitive operations;
  \item stopped-gradient approximation used only with an explicit correction or
    diagnostic label;
  \item not HMC-safe.
\end{itemize}
